% Generado por sistema/backend ReporteTex.java -- no editar a mano.
% Datos: CSV importado (titulados) / CSV importado (empleadores)
\subsubsection*{C14. Formación complementaria $\times$ situación laboral}

\textit{Paso 1. Describir el problema.} Factor de riesgo de desempleo: determinar si contar con formación complementaria (diplomado, especialidad, maestría o doctorado) se asocia con un menor riesgo de no tener trabajo.

\begin{itemize}[nosep]
  \item Variable independiente $X$: Formación complementaria.
  \item Variable dependiente $Y$: Situación laboral.
  \item Población: todos los titulados ($N = 148$). Criterio ARCU-SUR: 2.4.1 / 3.2.3. Nivel de significancia $\alpha = 0.05$.
  \item $H_0$: Las variables «Formación complementaria» (X) y «Situación laboral» (Y) son independientes.
  \item $H_1$: Existe dependencia significativa entre las variables «Formación complementaria» (X) y «Situación laboral» (Y).
\end{itemize}

\textit{Datos.} Primeros 10 de los 148 titulados que respondieron ambas preguntas (sin datos personales).

\begin{mitabla}{C14: datos de los primeros titulados}
\begin{tabular}{lll}
\toprule
Titulado & $X$: Formación complementaria & $Y$: Situación laboral \\
\midrule
T1 & Sí & Emprendimiento propio \\
T2 & Sí & Trabaja en una organización \\
T3 & Sí & Trabaja en una organización \\
T4 & Sí & Trabaja en una organización \\
T5 & Sí & Emprendimiento propio \\
T6 & Sí & Trabaja en una organización \\
T7 & Sí & Trabaja en una organización \\
T8 & Sí & Trabaja en una organización \\
T9 & Sí & No trabaja \\
T10 & Sí & Trabaja en una organización \\
\bottomrule
\end{tabular}
\end{mitabla}

\textit{Paso 2. Tabla de contingencia} (frecuencia observada $FO_{ij}$, $N = 148$). Cada celda cuenta a los titulados que respondieron a la vez la fila ($X$) y la columna ($Y$).

\begin{mitabla}{C14: tabla de contingencia, frecuencia observada $FO_{ij}$}
\begin{tabular}{L{3.8cm}C{3.00cm}C{3.00cm}C{3.00cm}|c}
\toprule
$X$: Formación complementaria \textbackslash{} $Y$: Situación laboral & Trabaja en una organización & Emprendimiento propio & No trabaja & Total \\
\midrule
Sí & 73 & 10 & 15 & 98 \\
No & 35 & 5 & 10 & 50 \\
\midrule
Total & 108 & 15 & 25 & $N = 148$ \\
\bottomrule
\end{tabular}
\end{mitabla}

\textbf{Mapa de calor.} La misma tabla de contingencia, pintada por color en vez de por número:

\begin{mitabla}{C14: mapa de calor de $FO_{ij}$}
\begin{tabular}{L{3.8cm}C{3.00cm}C{3.00cm}C{3.00cm}}
\toprule
$X$: Formación complementaria \textbackslash{} $Y$: Situación laboral & Trabaja en una organización & Emprendimiento propio & No trabaja \\
\midrule
Sí & \cellcolor{gray!75}73 & \cellcolor{gray!10}10 & \cellcolor{gray!15}15 \\
No & \cellcolor{gray!36}35 & \cellcolor{gray!5}5 & \cellcolor{gray!10}10 \\
\bottomrule
\end{tabular}
\end{mitabla}

\textit{Paso 3. Frecuencia esperada.} Es lo que habría en cada celda si $X$ e $Y$ fueran independientes. Por ejemplo, para la celda Sí y Trabaja en una organización:
\[
FE_{11} = \frac{(98)(108)}{148} = 71.51
\]

\begin{mitabla}{C14: frecuencia esperada $FE_{ij}$}
\begin{tabular}{L{3.8cm}C{3.00cm}C{3.00cm}C{3.00cm}}
\toprule
$X$: Formación complementaria \textbackslash{} $Y$: Situación laboral & Trabaja en una organización & Emprendimiento propio & No trabaja \\
\midrule
Sí & $98\cdot108/148 = 71.51$ & $98\cdot15/148 = 9.93$ & $98\cdot25/148 = 16.55$ \\
No & $50\cdot108/148 = 36.49$ & $50\cdot15/148 = 5.07$ & $50\cdot25/148 = 8.45$ \\
\bottomrule
\end{tabular}
\end{mitabla}

\textit{Paso 4. Estadístico chi-cuadrado.} En cada celda se calcula $\frac{(FO_{ij}-FE_{ij})^2}{FE_{ij}}$ y luego se suman todas.

\begin{mitabla}{C14: cálculo de $X_0^2$ celda por celda}
\begin{tabular}{L{3.2cm}C{3.20cm}C{3.20cm}C{3.20cm}|c}
\toprule
$X$: Formación complementaria \textbackslash{} $Y$: Situación laboral & Trabaja en una organización & Emprendimiento propio & No trabaja & Suma fila \\
\midrule
Sí & $\dfrac{(73-71.51)^2}{71.51}$\newline $= 0.0309$ & $\dfrac{(10-9.93)^2}{9.93}$\newline $= 0.0005$ & $\dfrac{(15-16.55)^2}{16.55}$\newline $= 0.1459$ & $0.1772$ \\[6pt]
No & $\dfrac{(35-36.49)^2}{36.49}$\newline $= 0.0606$ & $\dfrac{(5-5.07)^2}{5.07}$\newline $= 0.0009$ & $\dfrac{(10-8.45)^2}{8.45}$\newline $= 0.2859$ & $0.3474$ \\[6pt]
\bottomrule
\end{tabular}
\end{mitabla}

\[
X_0^2 = 0.1772 + 0.3474 = 0.5247
\]

\textit{Paso 5. Estadístico teórico.} Con $f = 2$ filas y $c = 3$ columnas: $gl = (2-1)(3-1) = 2$. De la tabla chi-cuadrada, con $\alpha = 0.05$: $X^2_{2,\,0.05} = 5.991$.

\textit{Paso 6. Relación entre estadísticos y conclusión.}

\begin{center}
\begin{tabular}{l}
$0.5247 < 5.991 \rightarrow$ verdad\\
Si ($X_0^2 < X^2_{(f-1)(c-1),\,\alpha}$)\\
\quad \fbox{``Independencia''}\\
Sino\\
\quad NO ``Independencia''\\
FinSi
\end{tabular}
\end{center}

No se rechaza $H_0$: con estos datos, «Situación laboral» no depende de «Formación complementaria».

\textit{Recomendación.} La situación laboral no depende de tener formación complementaria: no hay evidencia de que por sí sola reduzca el riesgo de desempleo en esta muestra.

